\chapter{Вещественная висячая вершина-селектор: запрет по условию первого порядка}

Проверка сложности выделила единственного чисто графового кандидата:
квадрат Краевского с висячей вершиной \(M_3(\mathbb R)\). Предполагалось,
что она устранит произвольную матрицу поколений на квадрате. Здесь это
утверждение проверяется до построения полной модели KO6.

\section{Поблочное правило первого порядка}

Неприводимый бимодуль имеет парную метку \(v=(i,j)\). Для блока Дирака
\(D_{ij,kl}\) условие первого порядка содержит
\begin{equation}
 (a_i-a_k)D_{ij,kl}(b_j-b_l),
\end{equation}
поэтому ненулевое ребро требует \(i=k\) или \(j=l\). Это стандартное
правило строки или столбца для диаграмм Краевского; см.
\path{arXiv:hep-th/9701081} и \path{arXiv:2207.04466}.

Уравнения распадаются по блокам оператора Дирака. Третья вершина,
присоединённая к одному концу, не входит в уравнение для другого блока.

\section{Размещения квадрата и висячей вершины}

Пусть
\begin{equation}
 A=(o,o),\quad B=(o,h),\quad C=(h,h),\quad D=(h,o)
\end{equation}
--- прямоугольник, удовлетворяющий условию первого порядка. Новое слагаемое
\(s=M_3(\mathbb R)\) допускает две минимальные вершины, смежные с \(A\):
\begin{equation}
 E_L=(s,o),\qquad E_R=(o,s).
\end{equation}
Оба размещения проходят правило строки или столбца. Транспонированные метки
возникают в античастичной половине, поэтому дополнение посредством \(J\) не
создаёт препятствия.

Однако прежнее ребро \(AB\) по-прежнему имеет вид
\begin{equation}
 D_{AB}=d_{AB}\otimes Y_{AB},
 \qquad Y_{AB}\in M_3(\mathbb R),
\end{equation}
если поколения являются пассивной кратностью. Новое уравнение первого
порядка ограничивает \(D_{AE}\), но не \(D_{AB}\).

\begin{theorem}[Невозможность висячего селектора]
Добавление висячей вершины \(M_3(\mathbb R)\) не уменьшает поколенческий
коммутант ни одного ранее существовавшего ребра цикла. Поэтому геометрия с
висячим селектором отвергается до проверок KO6, аномалий и гессиана.
\end{theorem}

\section{Проверка дополнением Шура}

Исключение висячей вершины изменяет лишь диагональный блок в точке её
присоединения:
\begin{equation}
 K_A(z)\mapsto
 K_A(z)-D_{AE}(K_E-z)^{-1}D_{EA}.
\end{equation}
Тождество для независимых блоков цикла отсюда не следует. Выбор
\(D_{AE}\), специально создающий нужный проектор ранга один, означал бы
введение селектора вручную.

\section{Упущенная возможность}

Алгебра поколений должна действовать непосредственно на ограничиваемом
ребре. Положим
\begin{equation}
 A=(G,o),\qquad B=(G,h),
 \qquad G=M_3(\mathbb R).
\end{equation}
Тогда ребро \(AB\) сплетает один и тот же неприводимый левый \(G\)-модуль,
а потому
\begin{equation}
 \operatorname{End}_{M_3(\mathbb R)}(\mathbb R^3)=\mathbb R I_3.
\end{equation}
Произвольная девятимерная матрица тем самым сводится к одному скаляру.

Это не висячая вершина: поколения превращаются из пассивной кратности в
координату алгебры четырёхвершинного прямоугольника. Предыдущая проверка
квадрата этот случай не рассматривала.

\begin{proposition}[Исправление реестра поиска]
Пятивершинный граф остаётся единственным чисто графовым одноцикловым
расширением, но его селекторное толкование неверно. Четырёхвершинный класс
был исчерпан только при пассивной кратности поколений. Прямоугольник с
активной алгеброй \(M_3(\mathbb R)\) является отдельным, ранее не
проверенным случаем.
\end{proposition}

\section{Условия прекращения направления}

Исправленный прямоугольник допускается ровно к одной проверке. Направление
закрывается, если внутренние автоморфизмы \(M_3(\mathbb R)\) вводят
нефиксированную поколенческую калибровочную связь или безмассовые бозоны;
условия вещественности и первого порядка требуют новых свободных матриц;
цикл теряет комплексную ориентацию; квартичный член снова имеет неверный
знак; идеал лишних 2-форм удаляет связанную кривизну; либо требуется
независимый центральный вес следа.

\section{Вывод}

\begin{protocol}
\begin{enumerate}
 \item графовый висячий селектор: \textbf{направление закрыто};
 \item размещение висячей вершины: \textbf{кинематически допустимо};
 \item ограничение ребра цикла посредством вершины: \textbf{не выполнено};
 \item прямоугольник с активной алгеброй поколений:
 \textbf{допущен к одной проверке};
 \item математическая родительская архитектура: \textbf{не получена};
 \item физическое замыкание: \textbf{не достигнуто}.
\end{enumerate}
\end{protocol}

Следующая проверка --- \path{version5_family_algebra_rectangle_gate}.
Машинный сертификат сохранён в
\path{s2t_v5_real_selector_leaf_ko6_gate_results.json}.